\section{Evaluation}
\label{sec:eval}

\subsection{The metric}

Our primary metric is the true-positive rate at a false-positive rate of $10^{-3}$, written
TPR@$10^{-3}$: we set the threshold so that one non-lens in a thousand is flagged and report the
fraction of real lenses recovered there. We interpolate the ROC curve to FPR${}=10^{-3}$ with
every intermediate point kept. Lens finding has been reported at this
point~\citep{parul2024}, and at lens rarity a false-positive rate near $10^{-3}$ or lower is
needed for a selected sample not to be dominated by non-lenses~\citep{discoveryengineE}. The
metric depends only on the ranking of scores, so it does not depend on the class ratio of the
test pool. We
report AUROC as a secondary check only: with about 26\,000 negatives against 98 positives it
can look healthy while the operating point is poor. The metric is discrete, with a step of
$1/98 \approx 0.010$ on folds 1--3 and $1/97$ on fold 4.

\subsection{Pooling}

Baseline A is one model scoring all five folds, so a single pooled ROC curve is valid for it.
Every adaptation arm trains a different model per fold, so no common threshold exists. We
report those arms as a fold mean with the across-fold standard deviation and always compare
fold mean with fold mean.

\subsection{Intervals and fold dependence}
\label{sec:intervals}

Each comparison is paired: two arms ran on the same folds, so we difference within fold and
report a paired $t$-interval on the four per-fold differences, with $\mathrm{df} = 3$. At
$n = 4$ a bootstrap has at most 35 distinct resamples (256 ordered draws), so we fixed the $t$-interval in
the pre-registration. The frozen comparison (Table~\ref{tab:frozen}) and the FreeMatch
unfrozen $-$ frozen contrast use paired bootstraps instead; we say so where they appear. An
interval that admits zero is not evidence that two arms are equal: it
means this measurement, at four folds and about 98 positives each, cannot separate them.

The four confirmatory folds are not independent. Any two folds share $\rvSharedTargetMin{}$
to $\rvSharedTargetMax{}$ of their $\rvTargetLensPool{}$ target lenses, and every fold uses
the same source set and the same test negatives. Variance estimates that ignore such overlap
between training sets can understate the true variance by a wide
margin~\citep{bengio2004, nadeau2003}, so our intervals are probably too narrow. We expect the
main design's contrasts against the control to survive this, because their intervals sit far
from zero. The narrower comparisons may not, and neither may the dilution experiment's BNM
contrast, whose interval clears zero by \dlGBnmLoAbsThree{}.

\subsection{The contrast set}

Table~\ref{tab:contrasts6} reports \contrastCount{} paired $t$-intervals and one bootstrap interval from the main design, beside
the four primary contrasts of the dilution experiment, at the 95\% level, with no
multiplicity correction; several share a control arm, so they are not independent. Four were
registered with the $\lambda=0$ control, before it ran but after the adaptation arms' results
were known (Section~\ref{sec:prereg}; marked $\dagger$): $\lambda=0$ against Baseline A,
and DANN, CORAL and FreeMatch each against $\lambda=0$, with $\lambda=0$ read at its budget of
4000. The same CORAL and FreeMatch contrasts with $\lambda=0$ read at 2000 and 1500 were
declared with them as sensitivity readouts (marked $\ddagger$). The other contrasts, including
BNM $-$ $\lambda=0$, were not among the four. The dilution experiment
(Section~\ref{sec:dilute-methods}) adds \dlNtWord{} paired $t$-intervals (\dlNtG{} for $G$,
\dlNtD{} for $D$) and \dlNbootWord{} bootstrap intervals, which compare its frozen stages with
those of the main design, again with no multiplicity correction. They share the $\lambda=0$
control with the contrasts above.

\subsection{The test pool, and the change we made after pre-registration}
\label{sec:testpool}

We drew the test negatives from the parent pool before we knew that rounded identifiers hid
duplicates. Matching by exact Euclid ID showed that \DOneLensNeg{} catalogued grade A and B
lenses and \DOneDeflNeg{} simulation host galaxies were among the \TestNegDrawn{} test
negatives. In each fold, $\rvOwnPosMin{}$ to $\rvOwnPosMax{}$ of the \DOneLensNeg{} were that fold's own
test positives, counted a second time as negatives. We wrote and hashed an exclusion rule
(\texttt{rule.json}, self-certifying digest \texttt{\rvRuleSha{}\ldots}) before recomputing anything, applied
it to the stored scores, and retrained nothing. The cleaned pool holds \TestNegDOne{}
negatives, and every number in the main text uses it.

The change raises every arm's TPR, by $\dPoolUpMin{}$ to $\dPoolUpMax{}$, and raises the
adaptation arms the most. On the original pool, the adaptation arms placed $\rvFsevenDaMin{}$
to $\rvFsevenDaMax{}$ of the \DOneLensNeg{} catalogued lenses among the $\rvAboveThrN{}$ negatives above
each fold's threshold, against $\rvFsevenRefMin{}$ to $\rvFsevenRefMax{}$ for Baseline A and
the control. The original pool therefore penalised the stronger models for ranking real
lenses highly. No contrast of the main design changes whether its interval includes zero
(Appendix~\ref{app:prereg_pool}). In the dilution experiment, whose analysis we pre-registered
on the cleaned pool, both forms of FreeMatch's contrast against the control exclude zero on the
original pool (Table~\ref{tab:dl-sensitivity}).

\rwDTwoWord{} of the test lenses also appear in training or validation as negatives: $\rvDTwoRandom{}$ as random galaxies, and one as a bare deflector
that also hosts simulated lenses. Fourteen of them were training negatives and entered the
loss labelled as non-lenses. The other six were validation negatives, which no loss uses.
Removing them changes each arm's fold-mean TPR by at most $\dPoolDTwoMax{}$
(Appendix~\ref{app:prereg_pool}).

\subsection{Secondary readouts and comparability to the benchmark}
\label{sec:eval-comparability}

We also report the completeness at 24\% purity, the operating point of the
benchmark~\citep{lines2025simstosky}. Purity depends on rarity, and our pool is
far denser in lenses than the sky, so we correct it to $N_N/N_P = 2200$, the ratio of objects
searched to lenses found in Q1, which~\citep{lines2025simstosky} adopts and we follow for
comparability, by
\begin{equation}
  \mathrm{purity} \;=\; \frac{\mathrm{TPR}}{\mathrm{TPR} + \mathrm{FPR} \cdot 2200}.
  \label{eq:purity}
\end{equation}
We take the completeness as the largest TPR among thresholds at which the corrected purity is
at least 24\%, since purity need not fall monotonically along the curve; a fold with no such
threshold is n/a, and an arm with an n/a fold has no fold mean and no interval. We
express completeness as a ratio to the benchmark's 0.50.

Lines et al.'s simulation-only model never sees a target image. Our Baseline A and $\lambda=0$
control are comparable to it in that respect, but our adaptation arms are not: they also see
a target built with catalogue labels, balanced 1:1 in the main design and diluted to 1:10 in
the second experiment (Section~\ref{sec:labels}). Their benchmark ratios compare a model given
such a target with one given none. The benchmark's
figure is a mean over ten initialisations with a spread of $\pm 0.10$ in purity, which our
intervals do not include, and its negatives are not matched to ours in depth
(Section~\ref{sec:splits}).
